\section{量词规则与等号}


\subsection{一个完整的量词例子}


目标是

\[
\neg\forall x\neg\forall y\varphi(x,y)
\vdash
\forall y\neg\forall x\neg\varphi(x,y).
\]

证明的关键是两次使用全称引入时都避开开放假设中的特征变量：

\begin{verbatim}
1 | ¬∀x¬∀y φ(x,y)       PR
2 |   ∀x¬φ(x,a)          AS
3 |     ∀yφ(b,y)         AS
4 |     ¬φ(b,a)          ∀E 2
5 |     φ(b,a)           ∀E 3
6 |     ⊥                →E 4,5
7 |   ¬∀yφ(b,y)          →I 3-6
8 |   ∀x¬∀yφ(x,y)        ∀I 7   (b 不在开放假设中)
9 |   ⊥                  →E 8,1
10| ¬∀x¬φ(x,a)           →I 2-9
11| ∀y¬∀x¬φ(x,y)         ∀I 10  (a 不在开放假设中)
\end{verbatim}


第 10 行是由假设 $[\forall x\neg\varphi(x,a)]$ 导出矛盾后用 $\to I$ 得到否定；不能把它误写成 RAA。

\subsection{等号规则}


等号推演可用下列规则概括：

\begin{itemize}
\item 自反：$t=t$；
\item 对称与传递：由 $t_1=t_2$ 得 $t_2=t_1$，由 $t_1=t_2,t_2=t_3$ 得 $t_1=t_3$；
\item 函数相容：逐个参数相等时，函数值相等；
\item 公式代换：在原子式或更复杂公式中，把等项替换为等项。
\end{itemize}


最后一条是语义中的 Leibniz 原则：等号两边在所有项和谓词中具有相同的解释。
